% =============================================================
%  Chapter — Short-Circuit Current
%  Figures in ./fig  (generated from short_circuit.py, cropped PDF).
%  Model: AC transient with decaying DC offset in an R-L circuit.
% =============================================================
\newcommand{\scfig}[1]{design/short-circuit/fig/#1}

\sectionslide{Short-Circuit Current}{AC transient with decaying DC offset}


% ---- Symbols ----
\begin{frame}{Symbols}
  \renewcommand{\arraystretch}{1.3}
  \begin{tabular}{@{}p{1.6cm}p{9.5cm}@{}}
    \textbf{Sym.} & \textbf{Description} \\
    \hline
    $u(t)$        & Source voltage, $u(t)=\hat{U}\sin(\omega t+\alpha)$ \\
    $\hat{U}$     & Peak source voltage, $\hat{U}=\sqrt{2}\,U_{\mathrm{rms}}$ \\
    $R,\;L$       & Circuit resistance and inductance \\
    $X$           & Reactance, $X=\omega L$ \\
    $Z$           & Impedance magnitude, $Z=\sqrt{R^{2}+X^{2}}$ \\
    $\varphi$     & Impedance angle, $\varphi=\arctan(X/R)$ \\
    $\alpha$      & Voltage phase at fault instant (switching angle) \\
    $\tau$        & DC time constant, $\tau=L/R=X/(\omega R)$ \\
    $i_p$         & Peak (making) current \\
    $\kappa$      & Asymmetry factor, $\kappa=i_p/(\sqrt{2}\,I_k'')$ \\
  \end{tabular}
\end{frame}

% ---- Abbreviations ----
\begin{frame}{Abbreviations}
  \renewcommand{\arraystretch}{1.3}
  \begin{tabular}{@{}p{1.7cm}p{9.5cm}@{}}
    \textbf{Abbr.} & \textbf{Meaning} \\
    \hline
    AC   & Alternating current \\
    DC   & Direct current \\
    RMS  & Root mean square \\
    SCC  & Short-circuit current \\
    X/R  & Reactance-to-resistance ratio \\
  \end{tabular}
\end{frame}


% ---- Physical model ----
\begin{frame}{The short-circuit current model}
  \begin{columns}[T,onlytextwidth]
    \begin{column}{0.5\textwidth}
      \begin{itemize}
        \setlength{\itemsep}{0.5em}
        \item A bolted fault on an $R$--$L$ branch fed from an ideal AC source
              $u(t)=\hat{U}\sin(\omega t+\alpha)$.
        \item The current is the sum of a \textbf{symmetrical AC} component and a
              decaying \textbf{DC (aperiodic)} component:
      \end{itemize}
      \vspace{0.6em}
      \[ i(t)=\frac{\hat U}{Z}\Big[\underbrace{\sin(\omega t+\alpha-\varphi)}_{\text{AC}}
              -\underbrace{\sin(\alpha-\varphi)\,e^{-t/\tau}}_{\text{DC}}\Big] \]
    \end{column}
    \begin{column}{0.5\textwidth}\centering
      \topoimg[width=\linewidth,height=0.62\textheight,keepaspectratio]{\scfig{waveform.pdf}}\\[0.2em]
      {\color{atGray}\footnotesize Worst-case fault at voltage zero ($\alpha=0$).}
    \end{column}
  \end{columns}
\end{frame}


% ---- Why the DC offset appears ----
\begin{frame}{Origin of the DC offset}
  \begin{itemize}
    \setlength{\itemsep}{0.6em}
    \item The current through an inductor \textbf{cannot change instantaneously}
          ($u_L=L\,\mathrm{d}i/\mathrm{d}t$), so at the fault instant $i(0)$ must
          equal its pre-fault value (zero for a fault from no load).
    \item The steady-state AC term alone would start at $\tfrac{\hat U}{Z}\sin(\alpha-\varphi)\neq 0$.
    \item The DC term is exactly the offset needed to satisfy the initial condition;
          it is the source-free (natural) response of the loop and decays with $\tau=L/R$.
    \item The resistance $R$ dissipates the associated energy $\Rightarrow$ exponential decay.
  \end{itemize}
\end{frame}


% ---- Key derived quantities ----
\begin{frame}{Key quantities}
  \renewcommand{\arraystretch}{1.4}
  \begin{columns}[c,onlytextwidth]
    \begin{column}{0.52\textwidth}
      \begin{tabularx}{\linewidth}{@{}lX@{}}
        Impedance      & $Z=\sqrt{R^{2}+X^{2}}$ \\
        Impedance angle& $\varphi=\arctan\!\big(X/R\big)$ \\
        AC peak        & $\hat{I}_{\mathrm{ac}}=\hat U/Z$ \\
        AC rms         & $I_k''=U_{\mathrm{rms}}/Z$ \\
        DC time const. & $\tau=L/R=X/(\omega R)$ \\
      \end{tabularx}
    \end{column}
    \begin{column}{0.44\textwidth}
      \begin{block}{Drivers of asymmetry}
        \begin{itemize}
          \item Switching angle $\alpha$ sets the \emph{size} of the DC offset.
          \item $X/R$ ratio sets \emph{how long} it persists ($\tau$).
        \end{itemize}
      \end{block}
    \end{column}
  \end{columns}
\end{frame}


% ---- Switching angle & peak making current ----
\begin{frame}{Switching angle and peak (making) current}
  \begin{itemize}
    \setlength{\itemsep}{0.5em}
    \item DC offset is \textbf{maximal} when $\alpha-\varphi=\pm\pi/2$
          (fault at a voltage zero for an inductive circuit) and \textbf{vanishes}
          when $\alpha-\varphi=0$ (fault at voltage peak).
    \item The first peak, the \textbf{peak making current}, reaches up to nearly
          twice the symmetrical AC peak for a high $X/R$:
  \end{itemize}
  \vspace{0.4em}
  \[ i_p=\kappa\,\sqrt{2}\,I_k'', \qquad
     \kappa=1.02+0.98\,e^{-3R/X}\ \ (\text{IEC }60909),\ \ 1\le\kappa\le 2 \]
  \vspace{0.2em}
  \begin{center}
    {\color{atGray}\footnotesize
    Example (230\,kV, $X/R=15$): $I_k''\approx 11.0$\,kA rms, $\kappa\approx 1.8$,
    $i_p\approx 28.3$\,kA.}
  \end{center}
\end{frame}


% ---- Same SCR, different X/R (Grid A vs Grid B) ----

\begin{frame}{Same fault level, different X/R}
\begin{tabular}{@{}l c c@{}}
        & Far from source & \makecell[c]{Close to source \\ (generator, \\ synchronous condenser)} \\
        & High $R$, lower $X$& Low $R$, high $X$\\
        \midrule
        SCR & $20$ & $20$ \\
        Sym. fault current $I_k''$& $25\,kA_{rms}$ & $25\,kA_{rms}$ \\
        X/R & $8$ & $40$ \\
         $\tau=\dfrac{X/R}{2 \cdot \pi \cdot f}$ & $\dfrac{8}{2\pi\cdot 50\,Hz}=25.5\,ms$ & $\dfrac{40}{2 \cdot \pi \cdot 50\,Hz}=127.3\,ms$
    % asdf & High $R$ &

  \end{tabular}

  \begin{block}{}
    \small Same SCR and same symmetrical fault current — but Grid DC offset of the grid close to source
    lasts $\sim\!5\times$ longer, giving a higher first current peak and more
    first-cycle stress on equipment.
  \end{block}  
\end{frame}


% ---- Comparison waveform ----
\begin{frame}{Same SCR, very different first-cycle stress}
  \begin{columns}[c,onlytextwidth]
    \begin{column}{0.6\textwidth}\centering
      \topoimg[width=\linewidth,height=0.66\textheight,keepaspectratio]{\scfig{scr_xr_compare.pdf}}
    \end{column}
    \begin{column}{0.38\textwidth}
      \begin{itemize}
        \setlength{\itemsep}{0.6em}
        \item Both curves share the same $I_{\mathrm{sym}}$ (same $\sqrt2\,I_{\mathrm{sym}}$ envelope).
        \item Higher X/R (Grid B) $\rightarrow$ slower DC decay $\rightarrow$
              higher peak making current.
        \item $i_p$: $\approx 60$ kA (X/R\,=\,8) vs $\approx 68$ kA (X/R\,=\,40).
      \end{itemize}
    \end{column}
  \end{columns}
\end{frame}


% ---- Why X/R rises near machines; SCR vs X/R ----
\begin{frame}{SCR vs.\ X/R}
  \begin{columns}[T,onlytextwidth]
    \begin{column}{0.5\textwidth}
      \textbf{\color{atBlack}Why X/R rises near machines}
      \begin{itemize}
        \setlength{\itemsep}{0.4em}
        \item Generators and synchronous condensers have high reactance $X$.
        \item Their winding resistance $R$ is very low.
        \item Closer electrical distance $\rightarrow$ more of their reactance is
              seen at the fault $\rightarrow$ X/R increases.
      \end{itemize}
    \end{column}
    \begin{column}{0.48\textwidth}
      \textbf{\color{atBlack}Two independent measures}
      \begin{itemize}
        \setlength{\itemsep}{0.5em}
        \item \textbf{SCR} tells you \emph{how much} fault current is available.
        \item \textbf{X/R} tells you \emph{how long} the asymmetry survives.
      \end{itemize}
      \vspace{0.4em}
      \begin{block}{}
        \small Two substations can have identical fault levels but place very
        different stresses on protection equipment.
      \end{block}
    \end{column}
  \end{columns}
\end{frame}


% ---- Energy view ----
\begin{frame}{Energy stored in the reactance}
  \begin{columns}[c,onlytextwidth]
    \begin{column}{0.56\textwidth}\centering
      \topoimg[width=\linewidth,height=0.66\textheight,keepaspectratio]{\scfig{energy.pdf}}
    \end{column}
    \begin{column}{0.42\textwidth}
      \begin{itemize}
        \setlength{\itemsep}{0.6em}
        \item $W_L=\tfrac{1}{2}L\,i(t)^2$ — energy scales with \emph{current squared}.
        \item Zero at $t=0$ (current cannot jump), then the asymmetric first
              half-cycle drives it well above the steady-state peak.
        \item The excess is dissipated in $R$ over $\tau$ as the offset decays.
      \end{itemize}
    \end{column}
  \end{columns}
\end{frame}


% ---- Takeaways ----
\begin{frame}{Summary}
  \begin{itemize}
    \setlength{\itemsep}{0.6em}
    \item SCC = symmetrical AC component $+$ decaying DC offset.
    \item The DC offset enforces continuity of inductor current; it decays with
          $\tau=L/R$ as $R$ dissipates the energy.
    \item $\alpha$ sets the magnitude of the offset (worst case at voltage zero);
          $X/R$ sets its persistence.
    \item Peak making current $i_p=\kappa\sqrt{2}I_k''$ sizes the mechanical and
          thermal withstand of breakers and busbars.
  \end{itemize}
\end{frame}
